% Phụ lục: chi tiết MMKG tiếng Việt (Khang).
% Nguồn số liệu: vietnamese-mmkg/{entities,triples,images,queries_auto}.jsonl, queries_hand.json, DATA_CARD.md,
% experiments/level3/vnmmkg/xlmr_h14/{a0_metrics,a0_explanations}.json. Điểm đồ thị ở Bảng tab:app-walk tính lại
% bằng EntityLinker + MkgwGraphScorer của scripts/eval_vnmmkg_transfer.py (load_kg), không cần encoder.
\appendices
\section{Chi tiết MMKG tiếng Việt}\label{app:vnmmkg}

Phụ lục này bổ sung cho mục~\ref{sec:vnmmkg}: quy trình dựng bộ dữ liệu, thống kê đầy đủ, vài mẫu dữ liệu thật và
một ví dụ chạy mô hình chuyển giao (kênh đếm cạnh KG với encoder XLM-R H/14) trên một câu, từng bước một.

\subsection{Quy trình dựng}

Hình~\ref{fig:app-pipeline} tóm tắt các bước của \texttt{python -m src.data.build\_vnmmkg}. Mọi lời gọi mạng
(Wikidata, Wikipedia, Commons) được cache trong \file{data/raw/vnmmkg/}, nên chạy lại từ cache cho ra đúng các file
đã commit. Hai bước có người duyệt (ảnh và triple) lưu quyết định trong \file{review_images.json} và
\file{review_triples.json}.

\begin{figure*}[t]
\centering
\footnotesize
\begin{tikzpicture}[
  box/.style={draw=cGray, rounded corners=2pt, align=center, text width=2.85cm, minimum height=1.35cm,
    inner sep=3pt, fill=white},
  hum/.style={box, draw=cTri, very thick},
  arr/.style={-{Stealth[length=2mm]}, cGray, thick},
  node distance=0.42cm]
\node[box] (sparql) {\textbf{Wikidata SPARQL}\\ \texttt{P17}=Việt Nam và có \texttt{P1435} hoặc thuộc 10 lớp công
  trình\\ 339 ứng viên};
\node[box, right=of sparql] (valid) {\textbf{Lọc}\\ có nhãn vi, bài viwiki, Commons category; bỏ hiện vật di
  động\\ 133 hợp lệ};
\node[box, right=of valid] (top) {\textbf{Xếp theo sitelink}\\ lấy 110 ứng viên đầu để duyệt ảnh};
\node[hum, right=of top] (img) {\textbf{Duyệt ảnh (tay)}\\ $\le$12 ứng viên/di tích; giữ 331, bỏ 541, quá
  hạn mức 144};
\node[box, right=of img] (lm) {\textbf{75 di tích}\\ 3--5 ảnh mỗi di tích; 35 di tích bị loại vì còn
  $<$3 ảnh};
\node[box, below=0.7cm of lm] (ctx) {\textbf{75 thực thể ngữ cảnh}\\ thêm theo bậc tới tổng 150 (bỏ 51 thực thể
  do giới hạn)};
\node[box, left=of ctx] (tri) {\textbf{Đề xuất triple}\\ Wikidata truthy, infobox viwiki (3 trường), 2 quy tắc
  (tỉnh, thời kỳ)};
\node[hum, left=of tri] (rev) {\textbf{Duyệt triple (tay)}\\ 107 triple infobox và suy diễn; bỏ 1};
\node[box, left=of rev] (kg) {\textbf{KG}: 325 triple, 13 quan hệ\\ Wikidata 257, infobox 26, suy diễn 42};
\node[box, left=of kg] (q) {\textbf{Truy vấn}\\ 197 tự động (câu đầu viwiki, che tên); 30 soạn tay, khoá
  trước khi chạy};
\draw[arr] (sparql) -- (valid);
\draw[arr] (valid) -- (top);
\draw[arr] (top) -- (img);
\draw[arr] (img) -- (lm);
\draw[arr] (lm) -- (ctx);
\draw[arr] (ctx) -- (tri);
\draw[arr] (tri) -- (rev);
\draw[arr] (rev) -- (kg);
\draw[arr] (kg) -- (q);
\end{tikzpicture}
\caption{Quy trình dựng MMKG tiếng Việt. Viền đỏ: bước duyệt tay. Số liệu lấy từ thống kê của bản build
(\texttt{DATA\_CARD.md}, mục 7).}
\label{fig:app-pipeline}
\end{figure*}

\subsection{Thống kê}

Bảng~\ref{tab:app-entities} chia 150 thực thể theo loại; Bảng~\ref{tab:app-relations} chia 325 triple theo quan
hệ và nguồn. 271 triple nối một di tích với một thực thể ngữ cảnh, 47 triple nối hai thực thể ngữ cảnh (ví dụ
Phật giáo \texttt{P138} Phật), 7 triple nối hai di tích (ví dụ Quần thể di tích Cố đô Huế \texttt{P527} Lăng
Khải Định). Mỗi di tích có trung bình 3,8 cạnh (trung vị 3, ít nhất 1, nhiều nhất 17 ở Quần thể di tích Cố đô
Huế); 7 di tích không có cạnh \texttt{P31} nên không được cộng điểm khi câu chỉ nhắc loại công trình. 9 thực thể
ngữ cảnh không có nhãn tiếng Việt (ví dụ ``Ramsar site'', ``Gothic Revival''), nên linker không bao giờ nối được
chúng.

\begin{table}[t]
\centering
\caption{150 thực thể theo loại.}
\label{tab:app-entities}
\footnotesize
\begin{tabular}{@{}p{1.9cm}rp{4.9cm}@{}}
\toprule
Loại & Số & Ví dụ \\
\midrule
Di tích & 75 & Chùa Một Cột, Nhà thờ Đức Bà Sài Gòn, Vịnh Hạ Long \\
Lớp & 17 & chùa, đền, nhà thờ Kitô giáo, nhà thờ chính tòa, lăng, tháp \\
Tỉnh, thành & 13 & Hà Nội, Thành phố Hồ Chí Minh, Huế, Khánh Hòa \\
Người & 7 & Lý Thái Tông, Thích Nhất Hạnh, Tự Đức \\
Tôn giáo & 5 & Phật giáo, Công giáo, đạo Cao Đài \\
Đối tượng thờ & 5 & Maria, Phật, Thánh Giuse \\
Danh hiệu & 5 & di sản thế giới, Bảo vật quốc gia của Việt Nam \\
Thời kỳ & 4 & nhà Lý, nhà Lê trung hưng, nhà Nguyễn, Pháp thuộc \\
Phong cách & 2 & Gothic Revival, Romanesque Revival \\
Khác & 17 & Hoàng thành Huế, Trung Bộ, Cột cờ Hà Nội \\
\bottomrule
\end{tabular}
\tabsrc{\file{vietnamese-mmkg/entities.jsonl}, trường \texttt{category}.}
\end{table}

\begin{table}[t]
\centering
\caption{325 triple theo quan hệ và nguồn (WD: Wikidata, IB: infobox, SD: suy diễn).}
\label{tab:app-relations}
\footnotesize
\setlength{\tabcolsep}{3pt}
\begin{tabular}{@{}llrrrr@{}}
\toprule
Quan hệ & Nghĩa & WD & IB & SD & Tổng \\
\midrule
\texttt{P31} & là một (lớp) & 86 & & & 86 \\
\texttt{P140} & tôn giáo & 40 & 17 & & 57 \\
\texttt{P131} & thuộc đơn vị hành chính & 29 & & 22 & 51 \\
\texttt{P1435} & danh hiệu di sản & 36 & & & 36 \\
\texttt{P527} & gồm có & 22 & & & 22 \\
\texttt{P2348} & thời kỳ & 1 & & 20 & 21 \\
\texttt{P361} & là một phần của & 20 & & & 20 \\
\texttt{P149} & phong cách kiến trúc & 9 & & & 9 \\
\texttt{P112} & người sáng lập & & 7 & & 7 \\
\texttt{P825} & thờ & 5 & 2 & & 7 \\
\texttt{P138} & đặt tên theo & 4 & & & 4 \\
\texttt{P276} & vị trí & 4 & & & 4 \\
\texttt{P186} & vật liệu & 1 & & & 1 \\
\midrule
Tổng & & 257 & 26 & 42 & 325 \\
\bottomrule
\end{tabular}
\tabsrc{\file{vietnamese-mmkg/triples.jsonl}, trường \texttt{r} và \texttt{source}.}
\end{table}

\begin{table}[t]
\centering
\caption{Ảnh và truy vấn.}
\label{tab:app-images-queries}
\footnotesize
\begin{tabular}{@{}p{4.6cm}r@{}}
\toprule
Thành phần & Số \\
\midrule
Ảnh đã duyệt & 331 \\
\quad di tích có 5 / 4 / 3 ảnh & 45 / 16 / 14 \\
\quad CC BY-SA / CC BY / Public Domain / CC0 & 163 / 83 / 57 / 28 \\
Truy vấn tự động & 197 \\
\quad di tích có 3 / 2 / 1 câu & 53 / 16 / 6 \\
\quad số từ (trung bình; ít nhất -- nhiều nhất) & 28,9; 7--90 \\
Truy vấn soạn tay (trực tiếp / gián tiếp) & 12 / 18 \\
\quad số từ (trung bình; ít nhất -- nhiều nhất) & 14,8; 7--30 \\
\quad câu gián tiếp có 1 / 2 / 3 / 4 đáp án & 15 / 1 / 1 / 1 \\
Chia A1 (val / test): di tích & 38 / 37 \\
\quad ảnh; câu tự động & 170 / 161; 94 / 103 \\
\bottomrule
\end{tabular}
\tabsrc{\file{vietnamese-mmkg/images.jsonl}, \file{queries_auto.jsonl}, \file{queries_hand.json},
\file{data/splits/vnmmkg_val.json}, \file{vnmmkg_test.json}.}
\end{table}

\subsection{Mẫu dữ liệu}

Hình~\ref{fig:app-samples} là bốn ảnh trong bộ, mỗi ảnh kèm các triple của di tích đó. Có thể thấy phần lớn
triple (tỉnh, thời kỳ, người sáng lập, danh hiệu) là điều không nhìn ra từ ảnh; đây là lý do kênh KG giúp được ở
câu gián tiếp. Bảng~\ref{tab:app-records} là một bản ghi thật của mỗi file, cùng lấy theo Chùa Một Cột.

\begin{figure*}[t]
\centering
\scriptsize
\setlength{\tabcolsep}{3pt}
\begin{tabular}{@{}cccc@{}}
\includegraphics[height=2.9cm]{vnmmkg/69.jpg} &
\includegraphics[height=2.9cm]{vnmmkg/41.jpg} &
\includegraphics[height=2.9cm]{vnmmkg/55.jpg} &
\includegraphics[height=2.9cm]{vnmmkg/122.jpg} \\[2pt]
\parbox[t]{4.1cm}{\textbf{(a) Chùa Một Cột} (ảnh 69)\\
\texttt{P31} chùa; \texttt{P131} Hà Nội; \texttt{P140} Phật giáo; \texttt{P2348} nhà Lý (SD, năm 1049 trong
infobox); \texttt{P112} Lý Thái Tông (IB).} &
\parbox[t]{4.1cm}{\textbf{(b) Nhà thờ Đức Bà Sài Gòn} (ảnh 41)\\
\texttt{P31} nhà thờ chính tòa, tiểu vương cung thánh đường; \texttt{P131} TP.~Hồ Chí Minh; \texttt{P140} Công
giáo; \texttt{P149} Romanesque Revival; \texttt{P2348} nhà Nguyễn (SD); \texttt{P825} Maria.} &
\parbox[t]{4.1cm}{\textbf{(c) Tháp Po Nagar} (ảnh 55)\\
\texttt{P31} đền; \texttt{P131} Khánh Hòa (SD). Chỉ có 2 triple: ít cạnh thì kênh KG ít giúp được.} &
\parbox[t]{4.1cm}{\textbf{(d) Lăng Khải Định} (ảnh 122)\\
\texttt{P31} lăng, tượng đài; \texttt{P131} Thừa Thiên - Huế (SD, tỉnh trước sáp nhập); \texttt{P1435} part of
UNESCO World Heritage Site; \texttt{P2348} Pháp thuộc (SD); \texttt{P361} Hoàng thành Huế.} \\
\end{tabular}
\caption{Bốn ảnh mẫu và triple của di tích tương ứng (IB: từ infobox, SD: suy diễn, còn lại từ Wikidata). Ảnh
Wikimedia Commons, đã thu nhỏ: (a) Cookie Nguyen, CC BY-SA 4.0; (b) Balon Greyjoy, CC0; (c) Ms Sarah Welch, CC0;
(d) dronepicr, CC BY 2.0.}
\label{fig:app-samples}
\end{figure*}

\begin{table*}[t]
\centering
\caption{Một bản ghi của mỗi file dữ liệu (rút gọn trường), cùng một di tích.}
\label{tab:app-records}
\footnotesize
\begin{tabular}{@{}p{2.7cm}p{14.3cm}@{}}
\toprule
File & Nội dung \\
\midrule
\file{entities.jsonl} & \texttt{qid} Q1186292, \texttt{kind} landmark, \texttt{label\_vi} ``Chùa Một Cột'',
\texttt{aliases\_vi} [``Chùa Hoa Sen'', ``Chùa Diên Hựu'', ``Diên Hựu'', ``Liên Hoa Đài'', ``Diên Hựu tự''],
\texttt{desc\_vi} ``công trình Phật giáo tại Hà Nội, Việt Nam'', \texttt{viwiki\_revid} 75560666,
\texttt{sitelinks} 17. \\
\file{triples.jsonl} & \texttt{h} Q1186292, \texttt{r} P2348, \texttt{t} Q878309 (nhà Lý), \texttt{source}
derived, \texttt{evidence} ``rule:period year 1049 from Thông tin chùa|founded=1049'', \texttt{reviewed} true. \\
\file{images.jsonl} & \texttt{image\_id} 69, \texttt{qid} Q1186292, \texttt{file\_title} ``File:Buddhist flags at
the One Piller Pagoda in 2016 (3).jpg'', 640$\times$360, \texttt{license} CC BY-SA 4.0, \texttt{artist} Cookie
Nguyen, \texttt{sha256} c1555ccf\ldots \\
\file{queries_auto.jsonl} & Câu gốc viwiki: ``\underline{Chùa Một Cột} có tên ban đầu là \underline{Liên Hoa
Đài} (\ldots) tức là Đài \underline{Hoa Sen} với lối kiến trúc độc đáo: một điện thờ đặt trên \underline{một
cột} trụ duy nhất.'' $\rightarrow$ truy vấn \texttt{Q1186292\_auto\_0}: ``có tên ban đầu là tức là Đài với lối
kiến trúc độc đáo: một điện thờ đặt trên trụ duy nhất.'' (phần gạch chân bị che; phần trong ngoặc bị bỏ). Cụm
``một cột'' bị che vì trùng tên sau khi bỏ tiền tố ``Chùa'', đúng hạn chế đã ghi trong data card. \\
\file{queries_hand.json} & \texttt{hand\_indirect\_01}, nhóm gián tiếp: ``Ngôi chùa ở Hà Nội được xây dựng dưới
thời nhà Lý''; \texttt{pattern} [(P31, chùa), (P131, Hà Nội), (P2348, nhà Lý)]; đáp án tính từ pattern trên
\file{triples.jsonl}: [Q1186292]. Câu trực tiếp \texttt{hand\_direct\_01} cùng đáp án: ``Ngôi điện gỗ nhỏ mái
ngói cong đặt trên một trụ đá duy nhất giữa hồ nước vuông, có cầu thang dẫn lên''. \\
\bottomrule
\end{tabular}
\end{table*}

\subsection{Ví dụ suy luận từng bước}\label{app:vn-walk}

Mô hình tốt nhất của cấp 2 (\texttt{kg\_triple}) chạy trên bộ tiếng Việt theo giao thức A0, không học gì thêm.
Với câu \texttt{hand\_indirect\_01} ``Ngôi chùa ở Hà Nội được xây dựng dưới thời nhà Lý'', các bước là:

\begin{enumerate}
\item \textbf{Truy hồi bằng encoder.} Câu và 331 ảnh được mã hoá bằng XLM-R large + ViT-H/14 (đóng băng); xếp ảnh
  theo cosine, giữ top-50.
\item \textbf{Liên kết thực thể.} \texttt{EntityLinker} đưa câu về chữ thường, tách âm tiết, so khớp chính xác
  với nhãn và alias tiếng Việt (tối đa 12 âm tiết). Ba thực thể được nối: ``hà nội'' $\rightarrow$ Q1858 (Hà
  Nội), ``nhà lý'' $\rightarrow$ Q878309 (nhà Lý), ``chùa'' $\rightarrow$ Q5393308 (lớp chùa).
\item \textbf{Đếm cạnh.} Với mỗi ảnh trong top-50, điểm đồ thị là số cạnh KG trực tiếp giữa di tích của ảnh và
  ba thực thể trên (Bảng~\ref{tab:app-walk}, Hình~\ref{fig:app-walk}).
\item \textbf{Kết hợp.} Điểm cuối $\alpha\, z(\mathrm{CLIP}) + (1-\alpha)\, z(\text{đồ thị})$ với $\alpha=0$ lấy
  từ MKG-W val, nên thứ tự do số cạnh quyết định; ảnh cùng số cạnh giữ thứ tự CLIP.
\item \textbf{Kết quả và giải thích.} Ảnh hạng 1 là ảnh 69 của Chùa Một Cột, đúng đáp án. Giải thích in ra là chính
  ba cạnh đã đếm: \texttt{P131}~$\rightarrow$~Hà Nội, \texttt{P2348}~$\rightarrow$~nhà Lý,
  \texttt{P31}~$\rightarrow$~chùa.
\end{enumerate}

\begin{table}[t]
\centering
\caption{Điểm đồ thị của câu ``Ngôi chùa ở Hà Nội được xây dựng dưới thời nhà Lý'' trên 75 di tích.}
\label{tab:app-walk}
\footnotesize
\setlength{\tabcolsep}{3pt}
\begin{tabular}{@{}lcp{3.6cm}@{}}
\toprule
Di tích & Cạnh & Cạnh khớp \\
\midrule
\textbf{Chùa Một Cột} (đáp án) & \textbf{3} & \texttt{P131} Hà Nội, \texttt{P2348} nhà Lý, \texttt{P31} chùa \\
Chùa Trấn Quốc & 2 & \texttt{P131} Hà Nội, \texttt{P31} chùa \\
Đền Lý Bát Đế & 2 & \texttt{P2348} nhà Lý, \texttt{P31} chùa \\
18 di tích khác & 1 & \texttt{P131} Hà Nội hoặc \texttt{P31} chùa \\
54 di tích còn lại & 0 & \\
\bottomrule
\end{tabular}
\tabsrc{linker và hàm đếm cạnh của \texttt{eval\_vnmmkg\_transfer.py}, tính trên mọi di tích; khi chạy thật chỉ
ảnh trong top-50 CLIP được chấm. Ảnh hạng 1: \texttt{a0\_explanations.json}.}
\end{table}

\begin{figure}[t]
\centering
\footnotesize
\begin{tikzpicture}[
  lm/.style={draw=cClip, very thick, rounded corners=2pt, fill=cClip!8, align=center, inner sep=3pt},
  ent/.style={draw=cTri, thick, rounded corners=2pt, fill=cTri!10, align=center, inner sep=3pt},
  oth/.style={draw=cGray!60, rounded corners=2pt, text=cGray, align=center, inner sep=2pt},
  hit/.style={cTri, very thick},
  miss/.style={cGray!60},
  lab/.style={font=\scriptsize\ttfamily, fill=white, inner sep=1pt}]
\node[ent] (hn) at (0,0) {Hà Nội};
\node[ent] (ly) at (2.6,0) {nhà Lý};
\node[ent] (chua) at (5.2,0) {chùa};
\node[lm] (mc) at (2.6,-2.6) {\textbf{Chùa Một Cột}\\ 3 cạnh};
\node[lm, draw=cGray, fill=white] (tq) at (0,-2.6) {Chùa Trấn Quốc\\ 2 cạnh};
\node[lm, draw=cGray, fill=white] (lbd) at (5.2,-2.6) {Đền Lý Bát Đế\\ 2 cạnh};
\node[oth] (pg) at (1.4,-3.9) {Phật giáo};
\node[oth] (ltt) at (3.8,-3.9) {Lý Thái Tông};
\draw[hit, opacity=.45] (tq) -- (hn);
\draw[hit, opacity=.45] (tq) -- (chua);
\draw[hit, opacity=.45] (lbd) -- (chua);
\draw[hit, opacity=.45] (lbd) -- (ly);
\draw[hit] (mc) -- node[lab, pos=.62] {P131} (hn);
\draw[hit] (mc) -- node[lab, pos=.5] {P2348} (ly);
\draw[hit] (mc) -- node[lab, pos=.62] {P31} (chua);
\draw[miss] (mc) -- node[lab, text=cGray] {P140} (pg);
\draw[miss] (mc) -- node[lab, text=cGray] {P112} (ltt);
\node[font=\scriptsize, text=cGray, align=center] at (2.6,0.6) {thực thể linker nối được từ câu};
\end{tikzpicture}
\caption{Đồ thị con quanh ví dụ. Đỏ: cạnh tới thực thể được nhắc trong câu (được đếm; cạnh nhạt là của hai ứng viên
kế tiếp); xám: cạnh có trong KG nhưng câu không nhắc. Chùa Một Cột hơn hai ứng viên kế tiếp một cạnh, đúng là cạnh không nhìn thấy trong ảnh.}
\label{fig:app-walk}
\end{figure}

\textbf{Một câu trực tiếp bị kênh KG làm sai.} Câu \texttt{hand\_direct\_02} ``Nhà thờ gạch đỏ có hai tháp chuông
nhọn cân xứng hai bên mặt tiền, ở giữa có cửa sổ hoa hồng tròn'' có đáp án Nhà thờ Đức Bà Sài Gòn. Linker nối
``nhà thờ'' vào lớp Q16970 và ``tháp'' vào lớp Q199451. Đáp án chỉ có \texttt{P31} tới lớp con nhà thờ chính tòa,
nên không được cạnh nào; ảnh hạng 1 là của Nhà thờ Tân Định (một nhà thờ thường, được 1 cạnh \texttt{P31} $\rightarrow$
nhà thờ Kitô giáo). Ở $\alpha=0$, một cạnh này thắng mọi chi tiết hình ảnh mà encoder bắt được. Đây là nguyên nhân
nhóm trực tiếp giảm R@1 từ 75,00 xuống 33,33 ở Bảng~\ref{tab:vn-a0}.

\textbf{Chạy lại.} \texttt{python scripts/eval\_vnmmkg\_transfer.py -{}-encoder xlmr\_h14} ghi
\file{a0_metrics.json}, \file{a1_metrics.json} và \file{a0_explanations.json} vào
\file{experiments/level3/vnmmkg/xlmr_h14/}; encoder đóng băng nên kết quả tất định.
